\documentclass[12pt]{report}
\usepackage[utf8]{inputenc}
\usepackage{fontspec}
\setmainfont{Times New Roman}
\renewcommand{\contentsname}{Mục lục}
\renewcommand{\listfigurename}{Danh sách hình vẽ}
\renewcommand{\listtablename}{Danh sách bảng}
\renewcommand{\figurename}{Hình}
\renewcommand{\tablename}{Bảng}
\renewcommand{\abstractname}{Tóm tắt}
\renewcommand{\bibname}{Tài liệu tham khảo}
\renewcommand{\chaptername}{Chương}
\usepackage{graphicx}
\usepackage{indentfirst}
\usepackage[colorlinks=true, linkcolor=blue, urlcolor=blue, citecolor=blue]{hyperref}
\usepackage{amsmath}
\usepackage{amssymb}
\usepackage{algorithm}
\usepackage{algpseudocode}
\usepackage{subcaption}
\usepackage[top=2.5cm,bottom=2.5cm,inner=1cm,outer=1cm]{geometry}
\usepackage{float}
\usepackage{placeins}
\usepackage{fancyhdr}
\usepackage{fancybox}
\usepackage{lastpage}
\usepackage{setspace}
\usepackage{enumitem}
\usepackage{tabularx}
\usepackage{booktabs}
\usepackage{listings}
\usepackage{xcolor}
\usepackage{array}
\usepackage{longtable}

% Code listing style từ template
\lstdefinestyle{javacode}{
    language=Java,
    basicstyle=\ttfamily\footnotesize,
    keywordstyle=\color{blue}\bfseries,
    commentstyle=\color{gray}\itshape,
    stringstyle=\color{red},
    numberstyle=\tiny\color{gray},
    numbers=left,
    numbersep=5pt,
    frame=single,
    frameround=tttt,
    backgroundcolor=\color{gray!5},
    showstringspaces=false,
    breaklines=true,
    breakatwhitespace=true,
    tabsize=4,
    captionpos=b
}

\usepackage{tikz}
\usetikzlibrary{shapes,arrows,positioning,fit,calc}

\lstset{
    language=Python,
    basicstyle=\ttfamily\footnotesize,
    keywordstyle=\color{blue}\bfseries,
    commentstyle=\color{gray}\itshape,
    stringstyle=\color{red!70!black},
    numberstyle=\tiny\color{gray},
    numbers=left,
    numbersep=8pt,
    frame=single,
    frameround=tttt,
    backgroundcolor=\color{gray!5},
    showstringspaces=false,
    breaklines=true,
    breakatwhitespace=true,
    tabsize=4,
    captionpos=b,
    xleftmargin=15pt,
    xrightmargin=5pt
}

% Setup Header and Footer
\pagestyle{fancy}
\fancyhf{}
\setlength{\headheight}{15pt}
\renewcommand{\headrulewidth}{0.4pt}
\fancyhead[L]{IT3160 - Nhóm 5 - Phân loại cảm xúc Amazon Reviews}
\fancyfoot[C]{\thepage}

% Có thể biên dịch khi ảnh nằm trong thư mục overleaf_assets ở repository,
% hoặc khi toàn bộ ảnh đã được tải trực tiếp lên thư mục gốc của Overleaf.
\graphicspath{{overleaf_assets/}{./}}
\newcommand{\safeincludegraphics}[2][]{%
    \IfFileExists{#2}{%
        \includegraphics[#1]{#2}%
    }{%
        \IfFileExists{overleaf_assets/#2}{%
            \includegraphics[#1]{overleaf_assets/#2}%
        }{%
            \fbox{\parbox[c][4cm][c]{0.86\linewidth}{%
                \centering\textbf{Thiếu tệp hình}\\[0.2cm]
                \texttt{\detokenize{#2}}\\[0.2cm]
                Hãy tải tệp này từ thư mục \texttt{overleaf\_assets} lên Overleaf.%
            }}%
        }%
    }%
}

\newcommand{\modelname}{BiLSTM + Attention}
\newcommand{\datasetname}{Amazon Reviews Dataset}
\newcommand{\codeword}[1]{\texttt{\detokenize{#1}}}

\title{\textbf{PHÂN LOẠI CẢM XÚC ĐA MỨC TRÊN AMAZON REVIEWS}\\
\large Mô hình trọng tâm BiLSTM kết hợp Additive Attention}
\author{Nhóm 5}
\date{Hà Nội, 2026}

\raggedbottom
\setlength{\parindent}{1.2cm}
\setlength{\parskip}{0.15em}
\onehalfspacing

\begin{document}
\begin{titlepage}
\begin{center}
{\fontsize{14pt}{1}\selectfont ĐẠI HỌC BÁCH KHOA HÀ NỘI} \\[0.2cm]
{\fontsize{14pt}{1}\selectfont \textbf{Trường Công nghệ Thông tin và Truyền thông}} \\[1.5cm]

{\fontsize{18pt}{1}\selectfont \textbf{BÁO CÁO BÀI TẬP LỚN}} \\[0.3cm]
{\fontsize{16pt}{1}\selectfont \textbf{Học phần: Nhập môn Trí tuệ nhân tạo}} \\[0.5cm]
Mã học phần: IT3160 \hspace{0.5cm} Mã lớp học: 167855 \\[2.5cm]

{\fontsize{16pt}{1}\selectfont \textbf{Đề tài:}} \\[0.5cm]
{\fontsize{20pt}{1.2}\selectfont \textbf{Phân loại cảm xúc đa mức trên Amazon Reviews\\ bằng mô hình BiLSTM kết hợp Attention}} \\[2.5cm]

\begin{tabular}{ll}
\textbf{Giảng viên hướng dẫn:} & PGS. TS. Ngô Văn Linh \\
\textbf{Nhóm sinh viên thực hiện:} & Trần Lê Ngọc Tâm \hspace{0.5cm} 202416346 \\
& Nguyễn Phong \hspace{1.15cm} 202400066 \\
& Nguyễn Ngọc Tuấn Anh \hspace{0.2cm} 202400029 \\
\end{tabular}

\vfill
{\fontsize{14pt}{1}\selectfont \textbf{Hà Nội, Tháng 06 Năm 2026}}
\end{center}
\end{titlepage}


\tableofcontents
\listoffigures
\listoftables



% ==================== ABSTRACT ====================
\chapter*{ABSTRACT}
\addcontentsline{toc}{chapter}{ABSTRACT}

Báo cáo trình bày bài toán phân loại đánh giá sản phẩm Amazon thành năm mức từ
1 đến 5 sao. Trọng tâm của nghiên cứu là kiến trúc \textbf{Bidirectional Long
Short-Term Memory kết hợp additive attention có mask padding}. Hai mô hình
SimpleRNN và Transformer Encoder nhỏ được cài đặt trên cùng pipeline, cùng cách
chia dữ liệu và cùng không gian từ vựng để đóng vai trò đối chứng, không phải là
đối tượng chính của đề tài.

Pipeline sử dụng 20.405 review hợp lệ, chia phân tầng 80\% cho huấn luyện và
20\% cho validation. Văn bản được chuẩn hóa, token hóa bằng từ vựng tối đa
50.000 từ, cắt hoặc đệm về 200 token và biểu diễn bằng GloVe 100 chiều có thể
tiếp tục cập nhật. Mô hình chính gồm Embedding, Dropout, BiLSTM 128 đơn vị trên
mỗi chiều, additive attention trên vector trạng thái 256 chiều, Dense 64 nút và
Softmax 5 lớp. Kết quả trên 4.081 mẫu validation cho thấy BiLSTM + Attention đạt
Accuracy 81,38\%, Macro F1 0,3727, Weighted F1 0,7549 và rating MAE 0,3328; tốt
hơn cả SimpleRNN và Transformer trong lần chạy được lưu.

Tuy vậy, kết quả cũng cho thấy một hạn chế nghiêm trọng: dữ liệu lệch mạnh về
lớp 1 sao và class weights đã bị tắt trong cả ba lần chạy. SimpleRNN và
Transformer không dự đoán đúng bất kỳ mẫu 2, 3 hoặc 4 sao nào. BiLSTM +
Attention có tạo dự đoán cho các lớp giữa nhưng recall vẫn rất thấp, lần lượt
chỉ 0,83\%, 3,64\% và 3,77\%. Do đó, Accuracy cao không đồng nghĩa với hệ thống
đã giải quyết tốt bài toán năm lớp. Báo cáo phân tích thẳng thắn hiện tượng mô
hình co cụm về hai cực 1 và 5 sao, các giới hạn của quy trình đánh giá hiện tại,
và đề xuất hướng cải thiện dựa trên cân bằng lớp, loss có tính thứ bậc, tập test
độc lập và phân tích attention.

\textbf{Keywords:} sentiment analysis, Amazon Reviews, BiLSTM, additive
attention, class imbalance, ordinal classification, Transformer.



% ==================== SECTION 1 ====================
\chapter{Giới thiệu và phát biểu bài toán}

\section{Bối cảnh}

Đánh giá sản phẩm là nguồn dữ liệu quan trọng đối với cả người mua, nhà bán
hàng và nền tảng thương mại điện tử. Khác với bài toán sentiment nhị phân chỉ
phân biệt tích cực và tiêu cực, thang 1--5 sao chứa nhiều mức độ: 1 sao thường
thể hiện bất mãn rõ rệt, 5 sao thể hiện hài lòng mạnh, trong khi 2--4 sao chứa
những nhận xét pha trộn, trung tính hoặc có điều kiện. Việc nhận diện đúng các
lớp giữa khó hơn đáng kể vì ranh giới ngữ nghĩa mờ và nhiều review đồng thời
chứa cả ưu điểm lẫn nhược điểm.

Ví dụ, câu ``The screen is excellent, but the battery is disappointing'' chứa
hai tín hiệu trái chiều. Một mô hình chỉ dựa vào một vài từ tích cực hoặc tiêu
cực rất dễ đẩy review về 1 hoặc 5 sao. Mô hình cần vừa đọc được ngữ cảnh theo
hai chiều, vừa xác định thành phần nào trong chuỗi có ảnh hưởng lớn hơn tới
đánh giá cuối cùng. Đây là động cơ lựa chọn BiLSTM kết hợp Attention.

\section{Mục tiêu}

Các mục tiêu cụ thể của project gồm:

\begin{enumerate}[label=\arabic*.]
    \item Xây dựng pipeline thống nhất để tiền xử lý Amazon Reviews và phân loại
    mỗi review vào một trong năm mức rating.
    \item Cài đặt chi tiết mô hình chính BiLSTM + additive attention, trong đó
    padding được mask trước phép softmax của attention.
    \item Dùng SimpleRNN và một Transformer Encoder nhỏ làm hai đường cơ sở để
    đánh giá lợi ích của ngữ cảnh hai chiều và attention.
    \item So sánh không chỉ bằng Accuracy mà còn bằng Macro Precision, Macro
    Recall, Macro F1, Weighted F1, F1 từng lớp, confusion matrix,
    thời gian huấn luyện và thời gian suy luận.
    \item Nhận diện các thất bại của mô hình thay vì chỉ báo cáo con số tổng hợp,
    đặc biệt là khả năng phân loại các lớp 2, 3 và 4 sao.
\end{enumerate}

\section{Phạm vi và nguyên tắc đọc kết quả}

Ba mô hình được huấn luyện từ cùng một split và cùng biểu diễn GloVe. Transformer
trong project là một encoder nhỏ huấn luyện cho bài toán hiện tại, \textbf{không
phải} BERT hay một language model pretrained quy mô lớn. Vì vậy, so sánh trong
báo cáo nhằm trả lời câu hỏi kiến trúc nào hoạt động tốt hơn trong đúng thiết
lập này, không nhằm kết luận phổ quát rằng BiLSTM luôn tốt hơn Transformer.

Ngoài ra, tập được gọi là validation vừa tham gia điều khiển EarlyStopping vừa
được dùng để báo cáo kết quả cuối cùng. Chưa có tập test độc lập. Các con số vì
thế nên được xem là kết quả thực nghiệm nội bộ của project, chưa phải ước lượng
cuối cùng về khả năng tổng quát hóa trên dữ liệu hoàn toàn chưa quan sát.

\section{Đóng góp chính}

Đóng góp cốt lõi của báo cáo là việc xây dựng và phân tích mô hình BiLSTM + Attention với các đặc điểm:
\begin{itemize}
    \item Sử dụng GloVe 100 chiều được tinh chỉnh riêng cho miền Amazon.
    \item Triển khai BiLSTM với không gian trạng thái 256 chiều.
    \item Tích hợp Additive Attention có khả năng học và hỗ trợ mask padding.
    \item Áp dụng Dropout (0,5) và L2 regularization (0,01) để chống overfitting.
    \item Đánh giá toàn diện qua nhiều độ đo (Accuracy, F1) và tài nguyên.
\end{itemize}

Bên cạnh đó, báo cáo thẳng thắn phân tích nguyên nhân mô hình thất bại ở các lớp thiểu số (2-4 sao), qua đó khẳng định Accuracy cao không đồng nghĩa với việc bài toán đã được giải quyết trọn vẹn.



% ==================== SECTION 2 ====================
\chapter{Dữ liệu và pipeline thực nghiệm}

\section{Nguồn dữ liệu}

Project sử dụng \datasetname{} trên Kaggle, handle
\codeword{dongrelaxman/amazon-reviews-dataset}. Hai cột được đọc là
\codeword{Review Text} và \codeword{Rating}. Sau khi bỏ giá trị thiếu, rating
không hợp lệ, review rỗng và bản ghi trùng theo cặp văn bản--rating, pipeline thu
được 20.405 mẫu hợp lệ.

Nhãn rating ban đầu thuộc tập $\{1,2,3,4,5\}$ và được chuyển về
$\{0,1,2,3,4\}$ khi huấn luyện để phù hợp với
\codeword{sparse_categorical_crossentropy}. Khi báo cáo, nhãn được đổi lại thành
1--5 sao.

\section{Phân bố lớp}

Tập validation có 4.081 mẫu và phản ánh rõ mức độ mất cân bằng của dữ liệu.
Bảng~\ref{tab:class-distribution} được lấy trực tiếp từ support trong
classification report của ba mô hình.

\begin{table}[H]
\centering
\caption{Phân bố nhãn trên tập validation}
\label{tab:class-distribution}
\begin{tabular}{lrr}
\toprule
\textbf{Lớp} & \textbf{Số mẫu} & \textbf{Tỷ lệ} \\
\midrule
1 sao & 2.591 & 63,49\% \\
2 sao & 241 & 5,91\% \\
3 sao & 165 & 4,04\% \\
4 sao & 239 & 5,86\% \\
5 sao & 845 & 20,71\% \\
\midrule
\textbf{Tổng} & \textbf{4.081} & \textbf{100\%} \\
\bottomrule
\end{tabular}
\end{table}

Nếu luôn dự đoán 1 sao, một baseline rất đơn giản đã đạt 63,49\% Accuracy.
Điều này cho thấy Accuracy phải được đọc cùng Macro F1 và recall từng lớp. Ba
lớp 2--4 sao cộng lại chiếm 15,80\%, nhưng lại là phần khó và có ý nghĩa nhất
để kiểm tra mô hình có thực sự học bài toán năm lớp hay chỉ tách hai cực.

\section{Làm sạch và biểu diễn văn bản}

Hàm \codeword{clean_text} thực hiện:

\begin{enumerate}
    \item chuyển văn bản về chữ thường;
    \item loại bỏ thẻ HTML;
    \item giữ các ký tự \codeword{a-z}, chữ số, khoảng trắng và dấu nháy đơn;
    \item gộp nhiều khoảng trắng liên tiếp;
    \item loại review rỗng sau làm sạch.
\end{enumerate}

Tokenizer Keras có giới hạn tối đa 50.000 từ và token
\codeword{<OOV>}. Tokenizer chỉ được fit trên tập train, vì vậy vocabulary của
validation không làm rò rỉ vào quá trình học. Tokenizer lưu trong artifact chứa
26.714 mục; cộng token padding tạo ra kích thước embedding thực tế
$V=26.715$. Mỗi chuỗi được post-padding hoặc post-truncation về
$T=200$ token.

Ma trận embedding được khởi tạo từ GloVe 6B 100d \cite{pennington2014glove}.
Các từ khớp GloVe nhận vector pretrained; các từ không khớp bắt đầu bằng vector
0. Lớp embedding đặt \codeword{trainable=True}, nên toàn bộ vector tiếp tục
được cập nhật trong quá trình huấn luyện. Token 0 được dùng cho padding và
\codeword{mask_zero=True}.

\section{Chia dữ liệu và chống leakage}

Dữ liệu được chia train/validation theo tỷ lệ 80/20, stratify theo nhãn và dùng
seed 42. Với 20.405 mẫu, kích thước tương ứng là 16.324 mẫu train và 4.081 mẫu
validation. Thứ tự đúng của pipeline là:

\[
\text{làm sạch}
\rightarrow \text{chia train/validation}
\rightarrow \text{fit tokenizer trên train}
\rightarrow \text{biến đổi cả hai tập}.
\]

Cách làm này tốt hơn việc fit tokenizer trước khi chia dữ liệu, bởi trường hợp
sau sẽ để vocabulary của validation ảnh hưởng gián tiếp tới biểu diễn đầu vào.

\section{Cấu hình huấn luyện chung}

\begin{table}[H]
\centering
\caption{Cấu hình thực tế lấy từ các tệp \codeword{metrics.json}}
\label{tab:common-config}
\begin{tabularx}{0.92\linewidth}{>{\bfseries}l X}
\toprule
Thành phần & Giá trị \\
\midrule
Số lớp & 5 \\
Độ dài chuỗi & 200 token \\
Giới hạn từ vựng & 50.000; kích thước thực tế 26.715 \\
Embedding & GloVe 6B, 100 chiều, có fine-tune \\
Batch size & 128 \\
Epoch tối đa & 10 \\
Optimizer & Adam, learning rate ban đầu 0,001 \\
Loss & Sparse categorical cross-entropy + L2 từ các lớp Dense \\
Dropout & 0,5 \\
L2 & 0,01 \\
Validation split & 0,2; stratified; seed 42 \\
EarlyStopping & Theo dõi \codeword{val_loss}, patience 3, khôi phục trọng số tốt nhất \\
ReduceLROnPlateau & Theo dõi \codeword{val_loss}, factor 0,5, patience 2, min LR $10^{-6}$ \\
Class weights của lần chạy & \textbf{Bật} (\codeword{use_class_weights=true}) \\
\bottomrule
\end{tabularx}
\end{table}



\section{Các độ đo đánh giá}

Với $K=5$ lớp, Accuracy được tính bởi
\[
\mathrm{Accuracy} = \frac{\sum_{i=1}^{N}\mathbb{I}(y_i=\hat{y}_i)}{N}.
\]

Precision, Recall và F1 cho lớp $k$ lần lượt là
\[
P_k=\frac{TP_k}{TP_k+FP_k}, \qquad
R_k=\frac{TP_k}{TP_k+FN_k}, \qquad
F1_k=\frac{2P_kR_k}{P_k+R_k}.
\]

Macro F1 xem mọi lớp có trọng số như nhau:
\[
\mathrm{MacroF1}=\frac{1}{K}\sum_{k=1}^{K}F1_k.
\]
Do đó, nếu mô hình bỏ qua lớp thiểu số, Macro F1 sẽ giảm mạnh ngay cả khi
Accuracy còn cao. Weighted F1 dùng support của từng lớp làm trọng số, nên gần
với chất lượng trên các lớp lớn hơn và dễ bị lớp 1 sao chi phối.

% ==================== SECTION 3 ====================
\chapter{Mô hình chính: BiLSTM kết hợp Additive Attention}

\section{Lý do lựa chọn kiến trúc}

SimpleRNN truyền một hidden state qua toàn chuỗi nên dễ suy giảm thông tin dài
hạn. LSTM bổ sung cell state và các cổng điều khiển để giảm vấn đề vanishing
gradient \cite{hochreiter1997lstm}. BiLSTM chạy hai LSTM theo hai hướng
\cite{schuster1997brnn}: chiều thuận đọc từ trái sang phải, chiều nghịch đọc từ
phải sang trái. Nhờ đó, biểu diễn tại một vị trí chứa cả ngữ cảnh trước và sau
từ đang xét.

Tuy nhiên, nếu chỉ lấy trạng thái cuối, toàn bộ review vẫn phải bị nén vào một
vector duy nhất. Additive attention \cite{bahdanau2015attention} cho phép mô
hình học mức quan trọng của từng timestep, rồi lấy tổng có trọng số. Đây là
điểm phù hợp với sentiment analysis: các từ hoặc cụm như ``not worth'', ``works
perfectly'', ``but'', ``disappointed'' không đóng góp như nhau.

\section{Sơ đồ kiến trúc tổng thể}

Hình~\ref{fig:bilstm-architecture} biểu diễn đúng luồng tensor trong file
\codeword{BiLSTMAttention/BiLSTM_Attention.py}. Sơ đồ được vẽ trực tiếp bằng
TikZ nên không cần thêm ảnh kiến trúc bên ngoài.

\begin{figure}[H]
\centering
\begin{tikzpicture}[
    node distance=0.52cm,
    every node/.style={font=\small},
    block/.style={draw, rounded corners, minimum width=10.4cm, minimum height=0.83cm,
                  align=center, fill=blue!5},
    core/.style={block, fill=green!9, draw=green!45!black},
    attn/.style={block, fill=orange!13, draw=orange!70!black},
    output/.style={block, fill=red!8, draw=red!55!black},
    arrow/.style={->, thick, >=stealth}
]
\node[block] (input) {Input token IDs: $(B,200)$};
\node[block, below=of input] (embed)
{Embedding GloVe trainable: $26.715 \times 100$\\Output: $(B,200,100)$; \codeword{mask_zero=True}};
\node[block, below=of embed] (drop1) {Dropout $p=0{,}5$};
\node[core, below=of drop1] (bilstm)
{Bidirectional LSTM: $128$ đơn vị chiều thuận + $128$ đơn vị chiều nghịch\\
\codeword{return_sequences=True}; Output: $(B,200,256)$};
\node[block, below=of bilstm] (drop2) {Dropout $p=0{,}5$};
\node[attn, below=of drop2] (attention)
{Masked additive attention\\
$W_a\in\mathbb{R}^{256\times256}$, $b_a\in\mathbb{R}^{256}$,
$v_a\in\mathbb{R}^{256\times1}$\\Context vector: $(B,256)$};
\node[core, below=of attention] (dense)
{Dense $256\rightarrow64$, ReLU, L2 $=0{,}01$};
\node[block, below=of dense] (drop3) {Dropout $p=0{,}5$};
\node[output, below=of drop3] (softmax)
{Dense $64\rightarrow5$, Softmax, L2 $=0{,}01$\\
$P(1\text{ sao}),\ldots,P(5\text{ sao})$};

\draw[arrow] (input) -- (embed);
\draw[arrow] (embed) -- (drop1);
\draw[arrow] (drop1) -- (bilstm);
\draw[arrow] (bilstm) -- (drop2);
\draw[arrow] (drop2) -- (attention);
\draw[arrow] (attention) -- (dense);
\draw[arrow] (dense) -- (drop3);
\draw[arrow] (drop3) -- (softmax);
\end{tikzpicture}
\caption{Kiến trúc chi tiết của mô hình BiLSTM + Attention được sử dụng}
\label{fig:bilstm-architecture}
\end{figure}

\section{Kích thước tensor và số tham số}

\begin{table}[H]
\centering
\caption{Chi tiết từng lớp của mô hình chính}
\label{tab:bilstm-layers}
\resizebox{\linewidth}{!}{
\begin{tabular}{llllr}
\toprule
\textbf{Thứ tự} & \textbf{Lớp} & \textbf{Cấu hình} & \textbf{Output shape} & \textbf{Tham số} \\
\midrule
1 & Input & 200 token integer & $(B,200)$ & 0 \\
2 & Embedding & $V=26.715$, $d=100$, GloVe, trainable & $(B,200,100)$ & 2.671.500 \\
3 & Dropout & $p=0,5$ & $(B,200,100)$ & 0 \\
4 & Bidirectional LSTM & 128 đơn vị mỗi chiều, return sequences & $(B,200,256)$ & 234.496 \\
5 & Dropout & $p=0,5$ & $(B,200,256)$ & 0 \\
6 & Additive Attention & masked, một context vector & $(B,256)$ & 66.048 \\
7 & Dense & 64, ReLU, L2 $=0,01$ & $(B,64)$ & 16.448 \\
8 & Dropout & $p=0,5$ & $(B,64)$ & 0 \\
9 & Dense & 5, Softmax, L2 $=0,01$ & $(B,5)$ & 325 \\
\midrule
\multicolumn{4}{r}{\textbf{Tổng}} & \textbf{2.988.817} \\
\bottomrule
\end{tabular}}
\end{table}

Riêng embedding chiếm 2.671.500 tham số, tương đương 89,38\% tổng tham số mô
hình. Vì thế, chênh lệch tổng số tham số giữa ba kiến trúc không lớn: phần lớn
ngân sách tham số nằm ở bảng từ, không phải khối tuần tự.

\section{Cơ chế LSTM}

Tại timestep $t$, với vector đầu vào $x_t$, hidden state $h_{t-1}$ và cell state
$c_{t-1}$, một LSTM chuẩn tính:

\begin{align}
i_t &= \sigma(W_i x_t + U_i h_{t-1} + b_i),\\
f_t &= \sigma(W_f x_t + U_f h_{t-1} + b_f),\\
o_t &= \sigma(W_o x_t + U_o h_{t-1} + b_o),\\
\tilde{c}_t &= \tanh(W_c x_t + U_c h_{t-1} + b_c),\\
c_t &= f_t \odot c_{t-1} + i_t \odot \tilde{c}_t,\\
h_t &= o_t \odot \tanh(c_t).
\end{align}

Forget gate $f_t$ đóng vai trò cốt lõi trong việc gạt bỏ các thông tin nhiễu từ quá khứ; input gate $i_t$ điều khiển mức độ dung nạp nội dung mới; trong khi output gate $o_t$ xác định phần cell state nào được phép phát ra. Điểm đột phá lớn nhất của kiến trúc LSTM nằm ở luồng truyền trạng thái $c_t$: việc sử dụng phép cộng ($\oplus$) thay vì phép nhân ma trận liên tiếp giúp gradient có thể trượt qua các cell mà không bị triệt tiêu (vanishing gradient), khắc phục hoàn toàn điểm yếu chí tử của SimpleRNN.

\section{Biểu diễn hai chiều}

Hai LSTM độc lập sinh:
\[
\overrightarrow{h_t}
=\mathrm{LSTM}_{fwd}(x_t,\overrightarrow{h_{t-1}}),
\qquad
\overleftarrow{h_t}
=\mathrm{LSTM}_{bwd}(x_t,\overleftarrow{h_{t+1}}).
\]

Output tại vị trí $t$ là phép nối
\[
h_t=[\overrightarrow{h_t};\overleftarrow{h_t}]
\in\mathbb{R}^{256}.
\]

Với tham số \codeword{return_sequences=True}, mô hình không chỉ nhả ra trạng thái cuối cùng mà trả về toàn bộ quỹ đạo ẩn $H=[h_1,\ldots,h_T]\in\mathbb{R}^{T\times256}$. Ma trận $H$ này là nền tảng tối quan trọng để cơ chế Attention phía sau có thể quan sát, so sánh và trích xuất thông tin thay vì bị ép nén toàn bộ ngữ cảnh một cách thô bạo vào một vector cuối cùng.

\section{Additive attention có mask}

Attention layer của project học ba tham số:
\[
W_a\in\mathbb{R}^{256\times256},\quad
b_a\in\mathbb{R}^{256},\quad
v_a\in\mathbb{R}^{256\times1}.
\]

Với mỗi hidden state $h_t$, score chưa chuẩn hóa được tính:
\begin{align}
u_t &= \tanh(h_tW_a+b_a),\\
e_t &= u_tv_a.
\end{align}

Đối với timestep padding, implementation thay $e_t$ bằng $-10^9$ trước
softmax. Sau đó:
\begin{align}
\alpha_t &= \frac{\exp(e_t)}{\sum_{j=1}^{T}\exp(e_j)},\\
c &= \sum_{t=1}^{T}\alpha_t h_t.
\end{align}

Do $\exp(-10^9)\approx 0$, padding gần như nhận trọng số bằng 0. Context vector
$c\in\mathbb{R}^{256}$ là biểu diễn duy nhất chuyển sang classifier. Tổng số
tham số attention là
\[
256^2+256+256=66.048.
\]

\begin{algorithm}[H]
\caption{Masked additive attention trong mô hình chính}
\label{alg:attention}
\begin{algorithmic}[1]
\Require $H\in\mathbb{R}^{B\times T\times256}$, mask
$M\in\{0,1\}^{B\times T}$
\State $E \gets \mathrm{squeeze}(\tanh(HW_a+b_a)v_a)$
\State $E \gets \mathrm{where}(M=1,E,-10^9)$
\State $A \gets \mathrm{softmax}(E,\mathrm{axis}=T)$
\State $C \gets \sum_{t=1}^{T} A_t H_t$
\State \Return $C$
\end{algorithmic}
\end{algorithm}

\section{Classifier và regularization}

Context vector đi qua Dense 64 nút với ReLU:
\[
z=\mathrm{ReLU}(W_cc+b_c).
\]
Lớp cuối sinh phân phối xác suất:
\[
\hat{y}=\mathrm{softmax}(W_oz+b_o).
\]

Hai kernel Dense có L2 regularization:
\[
\mathcal{L}
=\mathcal{L}_{CE}
+\lambda\left(\lVert W_c\rVert_2^2+\lVert W_o\rVert_2^2\right),
\qquad \lambda=0,01.
\]

Dropout 0,5 được đặt sau embedding, sau BiLSTM và sau Dense 64. Lưu ý đây là
\codeword{Dropout} thông thường trong mã nguồn, không phải
\codeword{SpatialDropout1D}. Việc ghi đúng chi tiết này quan trọng vì hai lớp có
cách loại bỏ phần tử khác nhau.

\section{Diễn biến huấn luyện}

\begin{figure}[H]
\centering
\safeincludegraphics[width=0.96\linewidth]{bilstm_attention_training_history.png}
\caption{Loss và Accuracy theo epoch của BiLSTM + Attention}
\label{fig:bilstm-history}
\end{figure}

Mô hình chính huấn luyện đủ 10 epoch. Train accuracy tăng từ 70,69\% lên
83,14\%; validation accuracy tăng nhanh và ổn định quanh 81\% từ epoch 3. Giá
trị validation loss thấp nhất xuất hiện ở epoch cuối, khoảng 0,6042. Learning
rate được giảm từ 0,001 xuống 0,0005 ở epoch cuối sau giai đoạn plateau.

Khoảng cách train--validation accuracy cuối cùng chỉ khoảng 1,76 điểm phần trăm,
nên biểu đồ không cho thấy overfitting tổng thể quá mạnh. Tuy nhiên, kết luận
này chỉ đúng với metric tổng hợp. Một mô hình có thể có đường validation đẹp
nhưng vẫn thất bại ở lớp thiểu số; confusion matrix cho thấy đúng hiện tượng đó.

\section{Confusion matrix của mô hình chính}

\begin{figure}[H]
\centering
\safeincludegraphics[width=0.72\linewidth]{bilstm_attention_confusion_matrix.png}
\caption{Confusion matrix của BiLSTM + Attention trên 4.081 mẫu validation}
\label{fig:bilstm-cm}
\end{figure}

Mô hình nhận diện rất tốt hai lớp cực trị:

\begin{itemize}
    \item lớp 1 sao: recall 97,72\%, F1 0,9167;
    \item lớp 5 sao: recall 91,36\%, F1 0,7988.
\end{itemize}

Ngược lại, các lớp giữa rất yếu:

\begin{itemize}
    \item 2 sao: chỉ đúng 2/241 mẫu; 218 mẫu bị đẩy về 1 sao;
    \item 3 sao: chỉ đúng 6/165 mẫu; 86 mẫu về 1 sao và 58 mẫu lên 5 sao;
    \item 4 sao: chỉ đúng 9/239 mẫu; 193 mẫu bị đẩy lên 5 sao.
\end{itemize}

Tổng cộng, tập validation có 645 mẫu thuộc 2--4 sao, nhưng mô hình chỉ tạo 60
dự đoán thuộc ba lớp này, tương đương 1,47\% toàn bộ dự đoán. Trong khi đó, tỷ
lệ thật của ba lớp giữa là 15,80\%. Nói cách khác, BiLSTM + Attention đã học
khá tốt trục sentiment tiêu cực--tích cực, nhưng chưa học tốt độ mịn năm mức.

\section{Điểm mạnh và giới hạn nội tại}

\textbf{Điểm mạnh} của kiến trúc là khả năng kết hợp ngữ cảnh hai chiều, xử lý
chuỗi có thứ tự, mask padding đúng và tạo một cơ chế chọn timestep có thể mở
rộng để giải thích dự đoán.

\textbf{Giới hạn} là toàn bộ chuỗi bị nén vào một context vector 256 chiều với
một bộ trọng số attention duy nhất. Review có nhiều khía cạnh như giá, độ bền,
giao hàng và chất lượng có thể cần nhiều ``góc nhìn'' attention. Hơn nữa, lớp
hiện tại chỉ trả về context vector, không trả $\alpha_t$, nên project chưa thể
vẽ heatmap từ quan trọng dù attention về lý thuyết hỗ trợ diễn giải. BiLSTM cũng
xử lý tuần tự và cần thấy toàn câu, không phù hợp bằng mô hình nhân quả nếu yêu
cầu streaming theo token.



% ==================== SECTION 4 ====================
\chapter{Hai mô hình đối chứng}

\section{SimpleRNN baseline}

SimpleRNN có luồng xử lý:
\[
\text{Embedding}
\rightarrow \text{Dropout}
\rightarrow \text{SimpleRNN}(128)
\rightarrow \text{Dense}(64)
\rightarrow \text{Dropout}
\rightarrow \text{Softmax}(5).
\]

Tại mỗi bước thời gian $t$, trạng thái ẩn (hidden state) được cập nhật thông qua phương trình đệ quy cơ bản:
\begin{equation}
h_t=\tanh(W_x x_t + W_h h_{t-1} + b).
\end{equation}
Trong đó, $W_x$ và $W_h$ là các ma trận trọng số biểu diễn quá trình tích hợp thông tin từ từ vựng hiện tại $x_t$ và ngữ cảnh quá khứ $h_{t-1}$.

Mô hình có 2.709.393 tham số, ít nhất trong ba kiến trúc. Việc bật cờ \codeword{unroll=True} giúp trải phẳng vòng lặp (unrolling) trên đồ thị tính toán theo chiều dài cố định 200, tối ưu hóa tốc độ nội suy tùy thuộc vào backend. Tuy nhiên, điểm yếu chí tử của SimpleRNN nằm ở chính phương trình đệ quy của nó: khi áp dụng thuật toán lan truyền ngược qua thời gian (BPTT - Backpropagation Through Time) trên một chuỗi dài 200 tokens, các đạo hàm phải nhân liên tiếp với ma trận $W_h$. Điều này dẫn đến sự bùng nổ (exploding) hoặc triệt tiêu đạo hàm (vanishing gradient) cực kỳ nghiêm trọng, khiến mạng gần như bất lực trong việc duy trì ngữ cảnh ở những từ đầu tiên trong câu.

Thực tế huấn luyện cho thấy đường validation của SimpleRNN dao động cực kỳ mạnh: độ chính xác (accuracy) có lúc tụt thẳng đứng xuống 49,72\% rồi bật lên trên 72\%. Hiện tượng bất định này là minh chứng rõ rệt nhất cho sự sụp đổ gradient của vanilla RNN khi phải vật lộn với nhiễu dữ liệu và Dropout 0,5. Trạng thái tốt nhất được lưu lại theo checkpoint có validation loss thấp nhất.

\section{Transformer Encoder baseline}

Transformer \cite{vaswani2017attention} là một bước ngoặt thay đổi hoàn toàn cách biểu diễn chuỗi bằng cách loại bỏ hoàn toàn cơ chế đệ quy, thay vào đó sử dụng Self-Attention để quan sát toàn cục. Luồng xử lý của khối Transformer như sau:
\[
\text{Token + Position Embedding}
\rightarrow \text{1 Transformer block}
\rightarrow \text{Masked Average Pooling}
\rightarrow \text{Dense}(64)
\rightarrow \text{Softmax}(5).
\]

Vì mô hình xử lý toàn bộ chuỗi song song nên khái niệm "thứ tự thời gian" bị xóa bỏ. Do đó, tín hiệu vị trí (Positional Encoding) bắt buộc phải được cộng trực tiếp vào Token Embedding. Lõi sức mạnh của khối encoder nằm ở phương trình \textbf{Scaled Dot-Product Attention}:
\begin{equation}
\text{Attention}(Q,K,V) = \text{softmax}\left(\frac{QK^T}{\sqrt{d_k}}\right)V
\end{equation}
Trong đó, Query ($Q$), Key ($K$) và Value ($V$) đều được ánh xạ tuyến tính từ cùng một nguồn biểu diễn đầu vào. Phép nhân ma trận $QK^T$ tạo ra một bản đồ tương quan giữa mọi cặp từ vựng trong câu bất chấp khoảng cách vật lý của chúng. Yếu tố tỷ lệ $\frac{1}{\sqrt{d_k}}$ đóng vai trò vô cùng quan trọng giúp giữ bình ổn phương sai của tích vô hướng, tránh việc hàm softmax bị đẩy vào vùng bão hòa gradient.

Mô hình triển khai \textbf{Multi-Head Attention} với 4 head độc lập (mỗi head có số chiều chiếu $\text{key\_dim} = 25$ để giữ nguyên tổng độ dài không gian là 100). Việc dùng nhiều head cho phép mạng linh hoạt nắm bắt đồng thời các mối quan hệ ngữ pháp và ngữ nghĩa khác nhau ở các không gian con (subspaces) riêng biệt. Lớp Feed-Forward Network (FFN) theo sau có cấu trúc mở rộng không gian $100\rightarrow256\rightarrow100$, tích hợp hàm kích hoạt GELU cùng cơ chế Residual Connection và Layer Normalization để chống suy biến mạng sâu. Tổng số lượng tham số là 2.790.645.

Về mặt lý thuyết, khả năng đối chiếu song song $O(N^2)$ của Transformer giúp nó vượt trội hoàn toàn về tốc độ chạy và khả năng phát hiện liên kết xa so với LSTM. Tuy nhiên, sức mạnh này đi kèm một cái giá: Transformer khát dữ liệu (data-hungry) và rất dễ học lệch cấu trúc nếu thiếu đi một pha tiền huấn luyện (pre-training) khổng lồ như BERT. Kết quả thực nghiệm phản ánh chính xác điểm yếu này: sau epoch 3, tập huấn luyện bắt đầu "học vẹt" (loss giảm mạnh xuống 0,5077) trong khi validation loss lại bật ngược lên 1,0037. Mạng đã bị quá khớp (overfitting) nghiêm trọng và EarlyStopping phải buộc dừng sớm ở epoch thứ 6.

\section{Điều kiện so sánh}

Ba mô hình dùng chung:

\begin{itemize}
    \item cùng tập train/validation, seed 42 và tokenizer fit trên train;
    \item cùng GloVe 100d, batch size 128 và độ dài chuỗi 200;
    \item cùng loss, Adam, Dropout 0,5, L2 0,01, EarlyStopping và scheduler;
    \item cùng trạng thái \codeword{use_class_weights=true}.
\end{itemize}

Sự thống nhất này làm phép so sánh kiến trúc có ý nghĩa hơn. Dù vậy, số epoch
thực tế khác nhau do EarlyStopping: Transformer dừng ở 6 epoch, còn RNN và
BiLSTM + Attention chạy đủ 10 epoch.

\section{Tóm tắt khác biệt kiến trúc}

\begin{table}[H]
\centering
\caption{So sánh định tính ba kiến trúc}
\label{tab:architecture-comparison}
\begin{tabularx}{\linewidth}{>{\bfseries}l X X X}
\toprule
Tiêu chí & SimpleRNN & BiLSTM + Attention & Transformer Encoder \\
\midrule
Vai trò & Baseline đơn giản & \textbf{Mô hình chính} & Baseline self-attention \\
Ngữ cảnh & Một chiều & Hai chiều & Hai chiều trong toàn chuỗi \\
Phụ thuộc xa & Yếu & Tốt hơn nhờ LSTM & Trực tiếp qua self-attention \\
Tổng hợp chuỗi & Hidden state cuối & Additive attention & Masked average pooling \\
Song song hóa theo token & Thấp & Thấp & Cao \\
Chi phí theo độ dài & Xấp xỉ tuyến tính & Xấp xỉ tuyến tính & Self-attention bậc hai \\
Khả năng giải thích hiện tại & Thấp & Có tiềm năng, nhưng chưa xuất weights & Attention map phức tạp \\
Rủi ro chính & Vanishing gradient & Chậm, bottleneck một context & Overfit, cần dữ liệu/pretraining \\
\bottomrule
\end{tabularx}
\end{table}



% ==================== SECTION 5 ====================
\chapter{Kết quả và so sánh chi tiết ba mô hình}

\section{Kết quả tổng hợp}

\begin{table}[H]
\centering
\caption{Kết quả trên cùng tập validation 4.081 mẫu}
\label{tab:main-results}
\resizebox{\linewidth}{!}{
\begin{tabular}{lrrrrr}
\toprule
\textbf{Mô hình} & \textbf{Accuracy} & \textbf{Macro P} & \textbf{Macro R}
& \textbf{Macro F1} & \textbf{Weighted F1} \\
\midrule
SimpleRNN & 0,7493 & 0,2847 & 0,3252 & 0,3029 & 0,6818 \\
\textbf{BiLSTM + Attention} & \textbf{0,8138} & \textbf{0,4722} &
\textbf{0,3946} & \textbf{0,3727} & \textbf{0,7549} \\
Transformer Encoder & 0,7988 & 0,3026 & 0,3735 & 0,3336 & 0,7328 \\
\bottomrule
\end{tabular}}
\end{table}

\begin{figure}[H]
\centering
\safeincludegraphics[width=0.98\linewidth]{quality_metrics.png}
\caption{So sánh các metric chất lượng tổng hợp}
\label{fig:quality-metrics}
\end{figure}

BiLSTM + Attention đứng đầu ở toàn bộ metric chất lượng:

\begin{itemize}
    \item hơn SimpleRNN 6,44 điểm phần trăm Accuracy;
    \item Macro F1 cao hơn SimpleRNN 23,05\% theo tỷ lệ tương đối;
    \item hơn Transformer 1,49 điểm phần trăm Accuracy;
    \item Macro F1 cao hơn Transformer 11,74\%;
\end{itemize}

Kết quả này ủng hộ lựa chọn mô hình chính. BiLSTM giúp giữ ngữ cảnh tốt hơn
SimpleRNN; attention tổng hợp chuỗi linh hoạt hơn việc chỉ dùng hidden state
cuối; đồng thời mô hình ổn định hơn Transformer nhỏ huấn luyện không pretrain.

Tuy nhiên, Macro F1 tốt nhất vẫn chỉ là 0,3727, thấp hơn nhiều so với Accuracy
0,8138. Chênh lệch này là bằng chứng định lượng rằng hiệu năng đang tập trung ở
các lớp lớn.

\section{So sánh F1 theo từng mức rating}

\begin{figure}[H]
\centering
\safeincludegraphics[width=0.94\linewidth]{per_class_f1.png}
\caption{F1-score theo từng lớp rating}
\label{fig:per-class-f1}
\end{figure}

\begin{table}[H]
\centering
\caption{F1-score từng lớp}
\label{tab:per-class-f1}
\begin{tabular}{lrrrrr}
\toprule
\textbf{Mô hình} & \textbf{1 sao} & \textbf{2 sao} & \textbf{3 sao}
& \textbf{4 sao} & \textbf{5 sao} \\
\midrule
SimpleRNN & 0,8607 & 0,0000 & 0,0000 & 0,0000 & 0,6538 \\
BiLSTM + Attention & \textbf{0,9167} & \textbf{0,0158} & \textbf{0,0649}
& \textbf{0,0674} & \textbf{0,7988} \\
Transformer & 0,9057 & 0,0000 & 0,0000 & 0,0000 & 0,7621 \\
\bottomrule
\end{tabular}
\end{table}

Biểu đồ cho thấy cả ba mô hình thực tế chủ yếu giải một bài toán gần giống phân
loại hai cực. SimpleRNN và Transformer có F1 bằng 0 ở cả ba lớp 2, 3 và 4 sao.
BiLSTM + Attention là mô hình duy nhất tạo ra một số true positive ở các lớp
này, nhưng mức cải thiện còn nhỏ về giá trị tuyệt đối.

F1 lớp 3 và 4 của mô hình chính cao hơn lớp 2, nhưng đều dưới 0,07. Vì vậy,
không nên mô tả rằng mô hình đã ``bắt tốt class thiểu số''. Mô tả đúng hơn là:
\textit{BiLSTM + Attention giảm mức độ sụp đổ lớp so với hai baseline, nhưng
chưa khắc phục được mất cân bằng}.

\section{So sánh confusion matrix}

\begin{figure}[H]
\centering
\begin{subfigure}[t]{0.31\linewidth}
    \centering
    \safeincludegraphics[width=\linewidth]{rnn_confusion_matrix.png}
    \caption{SimpleRNN}
\end{subfigure}
\hfill
\begin{subfigure}[t]{0.31\linewidth}
    \centering
    \safeincludegraphics[width=\linewidth]{bilstm_attention_confusion_matrix.png}
    \caption{BiLSTM + Attention}
\end{subfigure}
\hfill
\begin{subfigure}[t]{0.31\linewidth}
    \centering
    \safeincludegraphics[width=\linewidth]{transformer_confusion_matrix.png}
    \caption{Transformer}
\end{subfigure}
\caption{Confusion matrix của ba mô hình}
\label{fig:all-confusion}
\end{figure}

SimpleRNN và Transformer chỉ sử dụng hai cột dự đoán 1 sao và 5 sao. Điều này
không phải lỗi hiển thị: các cột 2--4 thực sự bằng 0. Transformer tốt hơn RNN
chủ yếu do phân tách hai cực chính xác hơn, không phải vì hiểu tốt hơn các mức
trung gian.

BiLSTM + Attention sử dụng đủ năm cột, nhưng phân bố dự đoán vẫn lệch:

\begin{table}[H]
\centering
\caption{Phân bố nhãn dự đoán của BiLSTM + Attention}
\label{tab:bilstm-predicted-distribution}
\begin{tabular}{lrr}
\toprule
\textbf{Nhãn dự đoán} & \textbf{Số mẫu} & \textbf{Tỷ lệ} \\
\midrule
1 sao & 2.933 & 71,87\% \\
2 sao & 12 & 0,29\% \\
3 sao & 20 & 0,49\% \\
4 sao & 28 & 0,69\% \\
5 sao & 1.088 & 26,66\% \\
\bottomrule
\end{tabular}
\end{table}

Các lỗi chủ yếu có cấu trúc hợp lý theo thứ tự rating: 2 sao bị kéo về 1 sao,
4 sao bị kéo về 5 sao, còn 3 sao chia sang cả hai cực. Điều này cho thấy mô
hình học được polarity nhưng decision boundary cho cường độ sentiment chưa đủ
tinh.

\section{Đường validation của ba mô hình}

\begin{figure}[H]
\centering
\safeincludegraphics[width=0.98\linewidth]{validation_curves.png}
\caption{Validation loss và validation accuracy của ba mô hình}
\label{fig:validation-curves}
\end{figure}

BiLSTM + Attention có đường validation ổn định nhất: loss giảm nhanh và duy trì
quanh 0,60--0,64. Transformer hội tụ nhanh trong ba epoch đầu nhưng validation
loss tăng sau đó, trong khi SimpleRNN dao động mạnh.

Cần lưu ý rằng mỗi đường chỉ có một lần chạy với seed 42. Không có dải độ lệch
chuẩn hoặc confidence interval. Một đường đẹp hơn trong một split chưa đủ để
kết luận độ ổn định thống kê. Thực nghiệm đầy đủ nên chạy nhiều seed và báo cáo
trung bình $\pm$ độ lệch chuẩn.

\begin{figure}[H]
\centering
\begin{subfigure}[t]{0.49\linewidth}
    \centering
    \safeincludegraphics[width=\linewidth]{rnn_training_history.png}
    \caption{SimpleRNN}
\end{subfigure}
\hfill
\begin{subfigure}[t]{0.49\linewidth}
    \centering
    \safeincludegraphics[width=\linewidth]{transformer_training_history.png}
    \caption{Transformer Encoder}
\end{subfigure}
\caption{Training history chi tiết của hai mô hình đối chứng}
\label{fig:baseline-history}
\end{figure}

\section{Chi phí tài nguyên}

\begin{table}[H]
\centering
\caption{Số tham số và thời gian ghi nhận trong artifact}
\label{tab:resources}
\resizebox{0.94\linewidth}{!}{
\begin{tabular}{lrrrr}
\toprule
\textbf{Mô hình} & \textbf{Tham số} & \textbf{Epoch}
& \textbf{Train (giây)} & \textbf{Inference (ms/mẫu)} \\
\midrule
SimpleRNN & \textbf{2.709.393} & 10 & 235,36 & 1,4516 \\
BiLSTM + Attention & 2.988.817 & 10 & 308,47 & 0,8103 \\
Transformer Encoder & 2.790.645 & 6 & \textbf{56,42} & \textbf{0,1859} \\
\bottomrule
\end{tabular}}
\end{table}

\begin{figure}[H]
\centering
\safeincludegraphics[width=0.98\linewidth]{resource_comparison.png}
\caption{So sánh số tham số, thời gian train và thời gian inference}
\label{fig:resources}
\end{figure}

Transformer nhanh nhất: thời gian train được ghi nhận nhanh hơn BiLSTM khoảng
5,47 lần và inference nhanh hơn khoảng 4,36 lần. Điều này phù hợp với khả năng
song song hóa self-attention. BiLSTM chậm nhất khi train do phải chạy LSTM theo
hai hướng.

Tuy nhiên, không nên xem các con số thời gian là hằng số của kiến trúc. Chúng
phụ thuộc GPU, phiên bản TensorFlow, batch size, warm-up, I/O và số epoch.
SimpleRNN có inference chậm hơn BiLSTM trong lần chạy này dù kiến trúc đơn giản
hơn; đây có thể là hệ quả của implementation \codeword{unroll=True}, kernel
backend và overhead đo. Kết luận an toàn là Transformer nhanh nhất
\textit{trong artifact hiện tại}, không phải SimpleRNN luôn chậm hơn BiLSTM.

\section{So sánh trực tiếp BiLSTM + Attention với SimpleRNN}

Hai mô hình cùng dùng recurrence, nhưng khác ở ba điểm:

\begin{enumerate}
    \item LSTM có cổng và cell state, SimpleRNN chỉ có một hidden state.
    \item BiLSTM đọc cả hai chiều, SimpleRNN chỉ đọc từ trái sang phải.
    \item Attention của mô hình chính tổng hợp toàn bộ chuỗi, SimpleRNN dùng
    output cuối.
\end{enumerate}

Chênh lệch 6,44 điểm Accuracy cho thấy ba thay đổi này
có giá trị thực nghiệm rõ ràng. Đặc biệt, chỉ BiLSTM + Attention tạo được true
positive cho lớp 2--4 sao. Đổi lại, mô hình chính có thêm khoảng 279 nghìn tham
số và mất thêm 73 giây huấn luyện trong lần chạy.

\section{So sánh trực tiếp BiLSTM + Attention với Transformer}

Transformer gần mô hình chính hơn về Accuracy nhưng yếu hơn về Macro Precision,
Macro F1. Transformer nhận diện hai cực tốt nhưng co cụm hoàn toàn, trong
khi BiLSTM còn duy trì một lượng nhỏ dự đoán ở lớp giữa. Với dataset hiện tại,
inductive bias tuần tự và cổng nhớ của LSTM có vẻ phù hợp hơn một Transformer
encoder nhỏ huấn luyện không pretrain.

Ngược lại, Transformer vượt trội về tốc độ. Nếu mục tiêu triển khai ưu tiên
throughput và bài toán có thể rút về nhị phân, Transformer nhỏ là lựa chọn đáng
cân nhắc. Nếu mục tiêu giữ đủ năm lớp và tối ưu chất lượng hiện có, BiLSTM +
Attention là lựa chọn hợp lý hơn, nhưng bắt buộc phải xử lý mất cân bằng.

\section{Mô hình được lựa chọn}

Dựa trên kết quả hiện tại, project lựa chọn \textbf{BiLSTM + Attention} làm mô
hình chính vì:

\begin{itemize}
    \item đứng đầu Accuracy, Macro Precision, Macro Recall, Macro F1 và Weighted
    F1;
    \item là mô hình duy nhất không sụp hoàn toàn về hai nhãn;
    \item learning curve ổn định hơn;
    \item kiến trúc phù hợp mục tiêu nghiên cứu về ngữ cảnh hai chiều và
    attention.
\end{itemize}

Lựa chọn này là tương đối trong ba mô hình, không che giấu rằng Macro F1 0,3727
vẫn thấp và lớp giữa chưa đạt chất lượng sử dụng thực tế.



% ==================== SECTION 6 ====================
\chapter{Vấn đề tồn tại, hướng cải thiện và kết luận}

\section{Vấn đề nghiêm trọng nhất: mất cân bằng và sụp đổ lớp}

Dữ liệu có 63,49\% lớp 1 sao nhưng class weights bị tắt. Cross-entropy tối ưu
trung bình trên mẫu, nên gradient từ lớp 1 sao áp đảo lớp 2--4. Một chiến lược
dự đoán chủ yếu 1 và 5 sao vẫn đạt Accuracy cao. Đây là nguyên nhân trực tiếp
nhất giải thích sự chênh lệch giữa Accuracy và Macro F1.

Mã nguồn đã có khả năng tính class weight bằng
\codeword{compute_class_weight(class_weight="balanced")}; bước đầu tiên cần làm
là bảo đảm \codeword{USE_CLASS_WEIGHTS=1} trong lần train chính thức và lưu lại
weights vào artifact. Sau đó phải so sánh:

\begin{itemize}
    \item không class weight;
    \item class weight cân bằng;
    \item oversampling lớp 2--4;
    \item focal loss hoặc class-balanced focal loss.
\end{itemize}

Metric chọn checkpoint cũng nên cân nhắc Macro F1 thay vì chỉ
\codeword{val_loss}. Nếu mục tiêu là năm lớp cân bằng, validation loss trung
bình có thể tiếp tục ưu tiên lớp lớn.

\section{Bài toán có thứ bậc nhưng loss đang xem nhãn là danh mục}

Rating 1--5 có quan hệ thứ tự, nhưng softmax cross-entropy coi năm lớp là năm
category độc lập. Trong loss, dự đoán 1 sao thành 2 sao và 1 sao thành 5 sao đều
chỉ là ``sai lớp''. Rating MAE chỉ được dùng sau huấn luyện, chưa tác động vào
gradient.

Các hướng phù hợp hơn gồm:

\begin{enumerate}
    \item ordinal regression với các ngưỡng tích lũy
    $P(y>1),\ldots,P(y>4)$;
    \item Earth Mover's Distance hoặc loss dựa trên khoảng cách rating;
    \item label smoothing có cấu trúc, phân bố một phần xác suất cho lớp lân cận.
\end{enumerate}

Những loss này khuyến khích mô hình hiểu 3 sao nằm giữa 1 và 5 sao, thay vì chỉ
học năm nhãn rời rạc.

\section{Hạn chế của quy trình đánh giá}

\begin{enumerate}
    \item \textbf{Không có test set độc lập.} Validation được dùng cho
    EarlyStopping và báo cáo, nên kết quả có thể lạc quan.
    \item \textbf{Chỉ một split và một seed.} Chưa biết mức biến thiên của kết
    quả.
    \item \textbf{Không có confidence interval.} Chênh lệch 1,49 điểm Accuracy
    giữa BiLSTM và Transformer chưa được kiểm định thống kê.
    \item \textbf{Thời gian chưa chuẩn hóa đầy đủ.} Chưa lưu loại GPU, phiên bản
    TensorFlow, số lần warm-up và peak memory.
    \item \textbf{Không có error analysis theo độ dài hoặc từ OOV.} Chưa biết
    mô hình thất bại nhiều hơn ở review dài, ngắn hay chứa từ miền chuyên biệt.
\end{enumerate}

Thiết kế thực nghiệm nên chuyển sang train/validation/test, ví dụ 70/15/15 hoặc
80/10/10, giữ test hoàn toàn đóng cho tới khi chốt mô hình. Nên chạy ít nhất
3--5 seed.

\section{Hạn chế tiền xử lý}

Regex hiện chỉ giữ chữ cái tiếng Anh \codeword{a-z}, chữ số và dấu nháy đơn.
Điều này phù hợp phần lớn review tiếng Anh nhưng có thể làm mất emoji, ký hiệu
nhấn mạnh, từ Unicode và thông tin dấu câu. Trong sentiment analysis, ``!!!'',
emoji hoặc viết hoa toàn bộ có thể là tín hiệu hữu ích.

Post-truncation tại 200 token bỏ toàn bộ phần cuối review dài. Nếu kết luận hoặc
rating rationale thường xuất hiện ở cuối, đây là mất mát đáng kể. Cần thống kê
phân bố độ dài và tỷ lệ review bị cắt trước khi chọn 200.

GloVe không được huấn luyện riêng cho miền Amazon và coverage không được lưu
trong \codeword{metrics.json}. Các từ sản phẩm, model code và lỗi chính tả có
thể bắt đầu bằng vector 0. Nên lưu số lượng hit/miss, thử FastText subword hoặc
embedding pretrained theo miền.

\section{Hạn chế riêng của BiLSTM + Attention}

\begin{itemize}
    \item Attention chỉ có một context vector, có thể trộn nhiều khía cạnh trái
    chiều trong review.
    \item Attention weights chưa được xuất ra, nên chưa có kiểm chứng định tính
    mô hình đang tập trung vào từ nào.
    \item L2 $=0,01$ và Dropout $=0,5$ là regularization khá mạnh; cần ablation
    để xác định có làm underfit lớp thiểu số hay không.
    \item BiLSTM tuần tự nên chậm hơn Transformer khi chuỗi dài.
    \item Không có recurrent dropout trong LSTM và không có normalization sau
    BiLSTM; đây là các lựa chọn có thể thử nhưng cần đánh giá thực nghiệm.
\end{itemize}

Để phân tích attention, layer có thể trả đồng thời $(c,\alpha)$. Sau đó vẽ màu
trên token cho các ca đúng, ca sai và review chứa mệnh đề đối lập. Cần thận
trọng: attention weight là tín hiệu diễn giải hữu ích, nhưng không tự động đồng
nghĩa với giải thích nhân quả.

\section{Hạn chế riêng của hai baseline}

\textbf{SimpleRNN} gặp dao động huấn luyện, không giữ phụ thuộc dài và bỏ hoàn
toàn ba lớp giữa. Cờ \codeword{unroll=True} với chuỗi 200 có thể tăng memory và
không đảm bảo nhanh hơn trên mọi phần cứng. Baseline vẫn hữu ích để chứng minh
lợi ích của LSTM, nhưng không phù hợp để triển khai bài toán hiện tại.

\textbf{Transformer} overfit sau khoảng ba epoch và cũng bỏ hoàn toàn lớp 2--4.
Một encoder block 100 chiều không pretrained khó học đồng thời ngữ pháp, ngữ
nghĩa miền và decision boundary năm lớp từ dataset nhỏ. Có thể thử pretraining
như DistilBERT, hoặc tăng dữ liệu và dùng warm-up/weight decay; tuy nhiên khi đó
chi phí và phạm vi so sánh sẽ thay đổi đáng kể.

\section{Lộ trình cải thiện đề xuất}

\begin{table}[H]
\centering
\caption{Thứ tự ưu tiên cải thiện project}
\label{tab:roadmap}
\begin{tabularx}{\linewidth}{c >{\bfseries}l X}
\toprule
\textbf{Ưu tiên} & \textbf{Thay đổi} & \textbf{Mục tiêu} \\
\midrule
1 & Bật class weights và retrain đủ ba model &
Kiểm tra trực tiếp nguyên nhân sụp lớp 2--4. \\
2 & Dùng Macro F1 làm metric chọn model &
Không để lớp 1 sao chi phối checkpoint. \\
3 & Tạo test set độc lập, chạy nhiều seed &
Có ước lượng tổng quát hóa và độ bất định đáng tin cậy. \\
4 & Thử ordinal loss &
Khai thác cấu trúc thứ tự 1--5 sao và giảm lỗi nhảy nhiều mức. \\
5 & Xuất attention weights &
Phân tích token quan trọng và phát hiện shortcut. \\
6 & Tuning dropout, L2, units, max length &
Giảm underfit/overfit bằng ablation có kiểm soát. \\
7 & Thử pretrained Transformer hoặc FastText &
Cải thiện biểu diễn từ và từ ngoài vocabulary. \\
\bottomrule
\end{tabularx}
\end{table}

\section{Kết luận}

Trong thiết lập hiện tại, BiLSTM + Attention là mô hình tốt nhất trong ba kiến
trúc. Mô hình đạt Accuracy 81,38\% và Macro F1 0,3727; vừa
chính xác hơn, vừa tạo lỗi rating nhỏ hơn SimpleRNN và Transformer. Kiến trúc
128 đơn vị mỗi chiều kết hợp additive attention có mask là lựa chọn hợp lý cho
review độ dài vừa, và learning curve cho thấy quá trình huấn luyện tương đối ổn
định.

Kết luận quan trọng hơn là project chưa giải quyết tốt phân loại năm mức.
Accuracy cao chủ yếu đến từ hai lớp 1 và 5 sao. Class weights bị tắt trong dữ
liệu lệch mạnh làm cả ba mô hình co cụm về hai cực; BiLSTM + Attention chỉ giảm
nhẹ, chưa loại bỏ hiện tượng này. Do đó, bước tiếp theo không nên chỉ tăng kích
thước mô hình. Ưu tiên đúng là sửa thiết kế thực nghiệm: cân bằng lớp, dùng
metric/loss phù hợp với lớp thiểu số và thứ tự rating, tách test set, chạy nhiều
seed và kiểm tra attention. Khi các bước đó được thực hiện, giá trị thật sự của
kiến trúc BiLSTM + Attention mới có thể được đánh giá công bằng.



% ==================== REFERENCES ====================
\chapter*{TÀI LIỆU THAM KHẢO}
\addcontentsline{toc}{chapter}{TÀI LIỆU THAM KHẢO}

\begin{thebibliography}{99}

\bibitem{hochreiter1997lstm}
S. Hochreiter and J. Schmidhuber,
``Long Short-Term Memory,''
\textit{Neural Computation}, vol. 9, no. 8, pp. 1735--1780, 1997.

\bibitem{schuster1997brnn}
M. Schuster and K. K. Paliwal,
``Bidirectional Recurrent Neural Networks,''
\textit{IEEE Transactions on Signal Processing}, vol. 45, no. 11,
pp. 2673--2681, 1997.

\bibitem{bahdanau2015attention}
D. Bahdanau, K. Cho, and Y. Bengio,
``Neural Machine Translation by Jointly Learning to Align and Translate,''
\textit{International Conference on Learning Representations}, 2015.

\bibitem{vaswani2017attention}
A. Vaswani et al.,
``Attention Is All You Need,''
\textit{Advances in Neural Information Processing Systems}, vol. 30, 2017.

\bibitem{pennington2014glove}
J. Pennington, R. Socher, and C. D. Manning,
``GloVe: Global Vectors for Word Representation,''
\textit{Proceedings of EMNLP}, pp. 1532--1543, 2014.

\bibitem{kingma2015adam}
D. P. Kingma and J. Ba,
``Adam: A Method for Stochastic Optimization,''
\textit{International Conference on Learning Representations}, 2015.

\bibitem{amazonDataset}
Amazon Reviews Dataset, Kaggle,
\url{https://www.kaggle.com/datasets/dongrelaxman/amazon-reviews-dataset},
truy cập trong quá trình thực hiện project.

\bibitem{tensorflow}
M. Abadi et al.,
``TensorFlow: Large-Scale Machine Learning on Heterogeneous Systems,''
2015. \url{https://www.tensorflow.org/}.

\end{thebibliography}

\newpage

\end{document}
